% part: intuitionistic-logic
% chapter: soundness-completeness
% section: completeness-thm

\documentclass[../../../include/open-logic-section]{subfiles}

\begin{document}

\olfileid[hi]{int}{sc}{cpl}

\olsection{पूर्णता प्रमेय}

\begin{thm}\ollabel{thm:completeness}
  यदि $\Gamma \Entails !A$, तो $\Gamma \Proves !A$।
\end{thm}

\begin{proof}
  हम प्रतिधनात्मक रूप सिद्ध करते हैं। मान लें $\Gamma \Proves/ !A$। तब
  \olref[lin]{lem:lindenbaum} से कोई अभाज्य समुच्चय
  ~$\Gamma^* \supseteq \Gamma$ ऐसा है कि $\Gamma^* \Proves/ !A$।
  \olref[mod]{defn:canonical-model} में परिभाषित $\Gamma^*$ के विहित
  प्रतिरूप~$\mModel{M(\Gamma^*)}$ पर विचार करें। किसी भी
  $!B \in \Gamma$ के लिए $\Gamma^* \Proves !B$। ध्यान दें कि
  $\Gamma^*(\emptyseq) = \Gamma^*$। सत्यता लेम्मा
  (\olref[tru]{lem:truth}) से प्रत्येक $!B \in \Gamma$ के लिए
  $\mSat{M(\Gamma^*)}{!B}[\emptyseq]$, और
  $\mSat/{M(\Gamma^*)}{!A}[\emptyseq]$। इससे
  $\Gamma \Entails/ !A$ सिद्ध होता है।
\end{proof}

\begin{prob}
  दिखाइए कि यदि $!A$ में केवल !!{propositional variable}, $\lor$ और~$\land$ आते
  हैं, तो $\Entails/ !A$। इससे निष्कर्ष निकालिए कि अंतःप्रज्ञावादी
  तर्कशास्त्र में $\lif$ को $\lor$ और~$\land$ से परिभाषित नहीं किया जा सकता।
\end{prob}
  
\begin{prob}
  पूर्णता प्रमेय का उपयोग करके सिद्ध कीजिए कि यदि
  $\Proves !A \lor !B$, तो $\Proves !A$ या $\Proves !B$। (संकेत: मान लें
  $\mSat/{M_1}{!A}$ और $\mSat/{M_2}{!B}$, और ऐसा नया प्रतिरूप~$\mModel{M}$
  बनाइए कि $\mSat/{M}{!A \lor !B}$।)
\end{prob}
  
\begin{prob} 
  दिखाइए कि यदि $\mModel{M}$ रैखिक क्रम वाला संबंधात्मक प्रतिरूप है, तो
  $\mSat{M}{(!A \lif !B)\lor(!B \lif !A)}$।
\end{prob}

\end{document}
